% Session 2 -- Leaving FLAIR: Procedure TPS-P-01 (AN-15: "beginning with
% leaving FLAIR, since the generator runs only from the InitechDOS prompt").
% Law 1 sources: os/flair/finder_menu.c (Special menu: ... Restart, Shut
% Down, Sleep, divider, Exit to InitechDOS, always enabled);
% os/flair/finder_windows.c finder_exit_dos_alert_text ("Are you sure you
% want to exit to InitechDOS?"); os/milton/kmain.c finder_exit_dos (buttons
% Exit and Cancel; Cancel leaves the desktop; Exit asks open applications to
% quit as Shut Down does, saves the desktop, re-initialises the console,
% prints the banner of spec/dos_banner.txt and starts COMMAND.COM as the
% permanent command processor on the data volume, drive A);
% os/milton/command.c (EXIT is ignored by the permanent processor; DATE and
% TIME prompts at startup without AUTOEXEC.BAT, initech-qskz; CLS; CD; DIR
% /W).  The screens are the journey route (screens/SCREENS, flair).
% GATES: test-flair-exit-to-dos test-flair-tps-journey test-pubs-screens test-pubs-tps
\Gates{test-flair-exit-to-dos, test-flair-tps-journey, test-pubs-screens, test-pubs-tps}
\chapter{Leaving FLAIR}

\SideHead{Introduction}The TPS Report Generator runs only from the
InitechDOS prompt.  Every run therefore begins at the FLAIR desktop with the
exit to InitechDOS.  In this session you complete the pre-run checklist,
leave the desktop, and make the directory that holds the program the
current directory.

\SideHead{Objectives}Upon completion of this session you will be able to:
\begin{enumerate}
\item Complete Checklist TPS-C-1.
\item Choose \textbf{Exit to InitechDOS} from the \textbf{Special} menu and
  confirm the exit.
\item Answer the date and time prompts and clear the screen.
\item Make \file{A:\bs APPS} the current directory and verify that the
  program files are present.
\end{enumerate}

\IsThisGoodForTheCompany

\section{Before the Run}
\index{pre-run checklist}
The pre-run checklist is completed at the desktop, before the exit.  Initial
each item as it is completed.  Retain the completed checklist with
\form{5}.

\begin{Checklist}{TPS-C-1}{Pre-Run}
\CheckItem{The current edition of this book is at the workstation.}
\CheckItem{\form{1} (Rev.\ 10/89) through \form{5} are at hand.}
\CheckItem{A copy of the memo about the new cover sheets is at hand, and its
  number is noted.}
\CheckItem{The report number for this run has been assigned and entered in
  \form{2}, column 1.}
\CheckItem{Work in every open application has been saved.}
\CheckItem{If the report now kept as \file{TPSREPT.BAK} must be retained
  longer, it has been copied as your supervisor directs.  The run deletes
  it (\proc{07}).}
\end{Checklist}

\section{The Exit to InitechDOS}
\index{Exit to InitechDOS}\index{Special menu}
\textbf{Exit to InitechDOS} is the last item of the \textbf{Special} menu,
in its own group below \textbf{Shut Down} and \textbf{Sleep}.  It is always
available.  Choosing it displays the alert \texttt{Are you sure you want to
exit to InitechDOS?} with two buttons, \textbf{Exit} and \textbf{Cancel}.

Clicking \textbf{Cancel}, or pressing \key{Esc}, removes the alert; the
desktop remains as it was.  Clicking \textbf{Exit} puts the desktop away as
\textbf{Shut Down} does: each open application is asked to quit, and if one
does not, the desktop remains.  The desktop is then saved, the screen
returns to the text display, and InitechDOS starts on the data volume,
drive A, with this banner:

\ReportLines{../../../spec/dos_banner.txt}{1}{2}

Because the data volume has no \file{AUTOEXEC.BAT}, InitechDOS asks for the
date and the time.  Pressing \key{Enter} at each prompt keeps the date and
the time shown.

\begin{Procedure}{TPS-P-01}{To Leave FLAIR for the InitechDOS Prompt}
\ProcField{Purpose}{To end the FLAIR session and make \file{A:\bs APPS} the
  current directory at the InitechDOS prompt.}
\ProcField{Frequency}{Each run of the TPS Report Generator.}
\ProcField{Performed by}{The operator.}
\ProcField{Prerequisites}{The FLAIR desktop is running with the shipped data
  volume as drive A.  Checklist TPS-C-1 is complete.}
\ProcField{Materials}{This book; Checklist TPS-C-1.}
\begin{ProcSteps}
\item Choose \textbf{Exit to InitechDOS} from the \textbf{Special} menu.
  \Response{Are you sure you want to exit to InitechDOS?}
\item Click \textbf{Exit}.  The desktop is put away and the InitechDOS
  banner is displayed, followed by the date prompt.
  \Response{Enter new date (mm-dd-yy):}
\item Press \key{Enter}.  The time prompt is displayed.
  \Response{Enter new time:}
\item Press \key{Enter}.  The prompt \texttt{A:\bs>} is displayed.
\item Type \type{CLS} and press \key{Enter}.  The screen is cleared and the
  prompt is displayed at the top.
\item Type \type{CD APPS} and press \key{Enter}.
  \Response{A:\bs APPS>}
\item Type \type{DIR TPS*.* /W} and press \key{Enter}.
  \OpsShot[marks={1=TPSREP.COM;2=A:\Bsl APPS>DIR}, label=fig:tpsdir, caption={The
  program files in \file{A:\bs APPS}.}]{tpsdir}
\end{ProcSteps}
\ProcField{Verification}{The listing names \file{TPSREP.COM},
  \file{TPSPARM.DAT} and \file{TPS.DBF} (callout 1).  The prompt is
  \texttt{A:\bs APPS>} (callout 2).}
\ProcField{Records retained}{Checklist TPS-C-1, with \form{5}.}
\ProcNote{If a file is missing from the listing, do not continue.  See your
  supervisor.}
\ProcNote{\cmd{EXIT} has no effect at this prompt.  The procedures in this
  book end at the InitechDOS prompt.}
\end{Procedure}

\begin{CalloutLegend}
\LegendRow{1}{File names}{The TPS Report Generator, its parameter record and
  the TPS table.}
\LegendRow{2}{Prompt}{The current directory is \file{A:\bs APPS}.  The
  program is run from this prompt.}
\end{CalloutLegend}

\section{The Current Directory}
\index{current directory}
The program reads its parameter record and its table from the current
directory.  If the current directory is \file{A:\bs}, InitechDOS does not
find the program at all:

\OpsShot[caption={The command typed before \proc{01}, step 6.}]{tpsroot}

If the command is typed with the directory name, the program starts, but
finds no parameter record in \file{A:\bs} and ends with return code 1:

\OpsShot[fromtext={A:\Bsl>APPS}]{tps001}

In both cases, type \type{CD APPS}, and run the program again from the
prompt \texttt{A:\bs APPS>}.

\section{Lab Project 1. Leaving FLAIR}
During this project you leave the desktop and verify the program files.
\begin{enumerate}
\item Complete Checklist TPS-C-1 at the desktop.
\item Carry out \proc{01}, steps 1 to 7.
\item Compare your screen with Figure~\ref{fig:tpsdir}.  The listing
  must name the same three files.
\item Write the time at which you completed the procedure here:
  \rule{30mm}{0.4pt}
\end{enumerate}

\section{Study Questions}
\begin{enumerate}
\item From which menu is \textbf{Exit to InitechDOS} chosen?
\item What happens when you click \textbf{Cancel} in the alert?
\item Why is \type{CD APPS} typed before the program is run?
\item What does InitechDOS display if \type{TPSREP TPS.DBF} is typed at the
  prompt \texttt{A:\bs>}?
\item What is done if a program file is missing from the listing?
\end{enumerate}
